\section{Background}
\label{sec:background}

\subsection{Liquid-Phase Transmission Electron Microscopy}
\label{sec:lptem}

A transmission electron microscope forms an image from the electrons that pass through
a thin sample, and resolves structure down to the sub-nanometre scale. A conventional
sample must be dry and in vacuum. In liquid-phase TEM (LPTEM) a thin layer of liquid is
sealed between two electron-transparent membranes, usually silicon nitride, so that
processes in solution can be filmed as they happen: nanoparticles diffusing, nucleating,
growing, dissolving and merging. These processes decide how catalysts, drug carriers and
new materials form, and LPTEM is one of very few ways to watch them directly.

The price is image quality. In bright-field imaging a metal particle appears dark
because it scatters electrons out of the beam, and the contrast follows Beer--Lambert
attenuation: the intensity behind a stack of materials falls exponentially with the
thickness of each material divided by its inelastic mean free path. The liquid and the
membranes attenuate the beam as well, so the thicker the liquid, the lower the
contrast. The dose must also be kept low, because the beam itself drives chemistry in the
liquid. At a few to a few tens of electrons per pixel per frame, shot noise, which
follows a Poisson distribution, dominates every frame, and small particles are barely
distinguishable from it. The detector's transfer function blurs the image further, and
the liquid cell often bulges, so the liquid thickness varies across the field of view.
Section~\ref{sec:simulator} builds each of these effects into the simulator.

\subsection{Physics of Nanoparticle Dynamics}
\label{sec:physics}

Nanoparticle dynamics in solution follow a few well-characterised laws, which is what
makes a quantitative test of a learned model possible.

\paragraph{Diffusion.} A particle undergoing Brownian motion in two dimensions has a
mean squared displacement $\langle |\Delta\mathbf{x}(t)|^2 \rangle = 4Dt$, where $D$ is
the diffusion coefficient. By the Stokes--Einstein relation $D$ falls with particle
radius, so larger particles diffuse more slowly. Departures from linear growth of the
MSD (anomalous diffusion) carry information about the particle's environment. The MSD
of a particle tracker is limited by the localisation error and by track length, so a
measured MSD should be compared with that of a perfect detector under the same
sampling.

\paragraph{Growth.} When a particle grows from monomers in solution, its radius follows a
power law $r \propto t^{\alpha}$ whose exponent identifies the rate-limiting step. If
growth is limited by diffusion of monomers to the surface, $\alpha = 1/3$; if it is limited
by the reaction at the surface, $\alpha = 1/2$. In Ostwald ripening, the
Lifshitz--Slyozov--Wagner theory~\cite{lifshitz1961,wagner1961}, small particles dissolve and
feed larger ones, the number of particles falls, the total volume is conserved, and the
mean radius coarsens as $t^{1/3}$. Estimating $\alpha$ from a noisy, short time series
of radii is the main quantitative test of this thesis.

\paragraph{Coalescence.} When two particles touch they can merge into one. A merge
conserves volume, so the new radius is $(r_1^3 + r_2^3)^{1/3}$, and it changes the
particle count and the size distribution abruptly. Unlike growth and diffusion, a merge
is an event, discrete and timed, which makes it a natural test of whether a model can
anticipate change.

\subsection{Joint-Embedding Predictive Architectures}
\label{sec:jepa}

Self-supervised video models learn from video alone by predicting one part of the input
from another. Generative models predict pixels. Joint-embedding predictive
architectures (JEPAs)~\cite{lecun2022path} instead predict in a learned representation
space: an encoder maps both the visible context and the hidden target to embeddings,
and a predictor is trained to output the target's embedding from the context's, with
the target encoder updated as an exponential moving average of the context encoder to
prevent collapse. Because the loss is computed on embeddings, the model need not spend
capacity on unpredictable detail, such as the exact pattern of shot noise, which a pixel
loss would force it to model or blur. I-JEPA applied the idea to images~\cite{ijepa},
V-JEPA to video~\cite{vjepa}.

V-JEPA~2~\cite{assran2025vjepa2} scales the approach to over a million hours of
internet video. Its ViT-L encoder, used in this thesis, splits a clip into tubelets of 2
frames by $16 \times 16$ pixels, embeds each as a token, and processes all tokens with a
vision transformer; it is trained to predict the embeddings of masked spatio-temporal
regions. Frozen V-JEPA representations have been shown to support intuitive physics
judgements on natural video, while pixel-prediction models stay near
chance~\cite{garrido2025intuitive}.

\paragraph{Latent world models.} A world model predicts how the state of an environment
evolves. A latent world model does so in the representation space of an encoder: it
encodes the past frames and rolls their latents forward with a learned predictor.
V-JEPA~2-AC trains an action-conditioned, block-causal predictor on top of the frozen
V-JEPA~2 encoder, with each frame encoded on its own, and uses it to plan robot actions
by minimising the latent distance to a goal image~\cite{assran2025vjepa2}. DINO-WM does
the same on frozen DINOv2 patch features~\cite{dinowm}. This thesis uses the same recipe
without actions: the encoder is frozen, the predictor forecasts latent frames, and the
forecasts are judged by reading physical observables off them.
